\subsection{Estructuras de Control}

\subsubsection{Regla: llaves obligatorias en todo bloque}

\textbf{Regla:} Llaves \texttt{\{\}} obligatorias en toda estructura de control, incluso con una única instrucción; regla \texttt{IDE0011} (Microsoft, 2023e).

\textbf{Código correcto:}
\begin{lstlisting}[style=csharpstyle]
if (diceRoll == MaximumFaceValue)
{
    GrantExtraTurn();
}
\end{lstlisting}

\textbf{Código incorrecto:}
\begin{lstlisting}[style=csharpstyle]
if (diceRoll == MaximumFaceValue)
    GrantExtraTurn();
\end{lstlisting}

\subsubsection{Regla: rama \texttt{default} obligatoria en \texttt{switch}}

\textbf{Regla:} Todo \texttt{switch} incluye una rama \texttt{default} explícita, y esa rama atiende solo los valores no contemplados: lanza \texttt{ArgumentOutOfRangeException} cuando el valor llegó como argumento, o \texttt{InvalidOperationException} cuando proviene del estado del objeto (McConnell, 2004). Prohibido usar la rama \texttt{default} para representar uno de los valores válidos; cada valor válido tiene su propio \texttt{case}.

\textbf{Descripción:} McConnell recomienda que la rama \texttt{default} detecte errores en lugar de esconder un caso legítimo (McConnell, 2004). Si la rama \texttt{default} devuelve uno de los resultados válidos, un valor agregado después al \texttt{enum} y olvidado en el \texttt{switch} recibe ese resultado en silencio, sin que ninguna prueba ni excepción lo revele.

\textbf{Excepción:} En código de presentación (colores, estilos, iconos), la rama \texttt{default} puede devolver un valor neutro cuando un valor desconocido no tiene consecuencias de negocio ni de seguridad. El valor neutro es un miembro con nombre, nunca un literal, y no coincide con el resultado de ningún otro \texttt{case}.

\textbf{Código correcto:}
\begin{lstlisting}[style=csharpstyle]
public TurnBadge SelectBadge(TurnState state)
{
    switch (state)
    {
        case TurnState.WaitingForRoll:
            return TurnBadge.Dice;
        case TurnState.Moving:
            return TurnBadge.Arrow;
        case TurnState.Finished:
            return TurnBadge.Flag;
        default:
            throw new ArgumentOutOfRangeException(nameof(state), state, "Unsupported turn state.");
    }
}

public Color SelectBorderColor(NoticeTone tone)
{
    switch (tone)
    {
        case NoticeTone.Error:
            return Palette.Error;
        case NoticeTone.Success:
            return Palette.Success;
        default:
            return Palette.Border;
    }
}
\end{lstlisting}

\textbf{Código incorrecto:}
\begin{lstlisting}[style=csharpstyle]
public TurnBadge SelectBadge(TurnState state)
{
    switch (state)
    {
        case TurnState.WaitingForRoll:
            return TurnBadge.Dice;
        case TurnState.Moving:
            return TurnBadge.Arrow;
        default:
            return TurnBadge.Flag;
    }
}
\end{lstlisting}
